\section{Operación privada por etiquetas}\label{sec:anexo-despliegue}

La guía \rutarepo{docs/deployment.md} especifica secretos, conexión, preparación
del host y publicación. El operador etiqueta una revisión verificada con una
versión como \texttt{v1.0.0}, sin mover etiquetas existentes. El artefacto y
\texttt{/version.json} identifican versión y commit.

La instalación objetivo utiliza VPS3, cuenta \texttt{qadra} y raíz
\texttt{/home/qadra/qadra}; el frontend usa un túnel SSH por el puerto 22022.
\texttt{journalctl --user} muestra registros; el controlador permite consultar el estado y recuperar la versión anterior,
preservando datos actuales. Una restauración SQL exige los controles de identidad
de \rutarepo{docs/database-operations.md}. Certificados, CRL y respaldos requieren
mantenimiento explícito. Cada versión incluye qpdf y FFmpeg/ffprobe fijados;
la preparación de servicios y la compilación del paquete conservan evidencia
separada de la primera activación y de una restauración real.


El manifiesto prospectivo se especifica en
\rutarepo{docs/backup-capture-manifest.md}. Debe ligarse a un recibo externo
confiable; conservar datos y huellas juntos no prueba autenticidad. El registro
\texttt{maintenance/restore/active.json} queda fuera del archivo de configuración
y mantiene cerrado el ingreso normal cuando está presente. La guía
\rutarepo{docs/deployment-restore-fence.md} distingue esa consulta de admisión de
la restauración todavía por componer: no ofrece un comando para borrar el
registro, ignorarlo o reabrir automáticamente. Actualizar el paquete de la
aplicación tampoco instala controladores ni cambia unidades de servicio.
